% Created 2025-12-27 sáb 14:03
% Intended LaTeX compiler: pdflatex
\documentclass[11pt]{article}
\usepackage[utf8]{inputenc}
\usepackage[T1]{fontenc}
\usepackage{graphicx}
\usepackage{longtable}
\usepackage{wrapfig}
\usepackage{rotating}
\usepackage[normalem]{ulem}
\usepackage{amsmath}
\usepackage{amssymb}
\usepackage{capt-of}
\usepackage{hyperref}
\usepackage{minted}

\usepackage{parskip}
\usepackage[utf8]{inputenc} %% For unicode chars
\usepackage{placeins}
\usepackage[margin=2.50cm]{geometry}
\usepackage[T1]{fontenc}
\usepackage{mathpazo}
\linespread{1.05}
\usepackage[scaled]{helvet}
\usepackage{courier}
\hypersetup{colorlinks=true,linkcolor=blue}
\RequirePackage{fancyvrb}
\DefineVerbatimEnvironment{verbatim}{Verbatim}{fontsize=\small,formatcom = {\color[rgb]{0.5,0,0}}}
\usepackage{mdframed}
\BeforeBeginEnvironment{minted}{\begin{mdframed}}
\AfterEndEnvironment{minted}{\end{mdframed}}
\setcounter{secnumdepth}{3}
\date{}
\title{}
\hypersetup{
 pdfauthor={},
 pdftitle={},
 pdfkeywords={},
 pdfsubject={},
 pdfcreator={Emacs 27.1 (Org mode 9.6.18)}, 
 pdflang={Spanish}}
\begin{document}

\setcounter{tocdepth}{3}
\tableofcontents

\setlength\parindent{10pt}

\section{Categories management}
\label{sec:org0aa5a93}
 Let's create a \texttt{Categories} class in PHP, which will manage
the CRUD operations for the categories in the To-Do application. Along
with the PHP class, I'll provide an example of HTML forms for adding and
viewing categories.

\subsection{PHP \texttt{Categories} Class}
\label{sec:orge232c6b}
\begin{minted}[]{php}
class Categories {
    private $pdo;

    public function __construct($pdo) {
        $this->pdo = $pdo;
    }

    public function create($name) {
        $sql = "INSERT INTO categories (name) VALUES (?)";
        $stmt = $this->pdo->prepare($sql);
        $stmt->execute([$name]);
    }

    public function getAll() {
        $sql = "SELECT * FROM categories";
        $stmt = $this->pdo->query($sql);
        return $stmt->fetchAll();
    }

    // Additional methods for updating and deleting categories can be added here
}
\end{minted}

\subsection{Explanation}
\label{sec:org60f192e}
\begin{itemize}
\item \textbf{Constructor}: Initializes the database connection.
\item \textbf{Create Category}: Adds a new category to the database.
\item \textbf{Get All Categories}: Retrieves all categories from the database.
\end{itemize}

\subsection{HTML Forms}
\label{sec:org9f48222}
\subsubsection{Form to Add a New Category (\texttt{add\_category.html})}
\label{sec:org6e2a987}
\begin{minted}[]{php}
<!DOCTYPE html>
<html>
<head>
    <title>Add New Category</title>
</head>
<body>
    <h1>Add New Category</h1>
    <form action="add_category.php" method="post">
        Category Name: <input type="text" name="name" required><br>
        <input type="submit" value="Add Category">
    </form>
</body>
</html>
\end{minted}

\subsubsection{PHP Script to Handle the Form (\texttt{add\_category.php})}
\label{sec:org85e8007}
\begin{minted}[]{php}
<?php
require_once 'Categories.php';
require_once 'db_connect.php'; // PDO connection

$categories = new Categories($pdo);

if ($_SERVER["REQUEST_METHOD"] == "POST") {
    $categories->create($_POST['name']);
    echo "Category added successfully.";
    // Redirect or show a success message
}
\end{minted}

\subsubsection{Displaying Categories (\texttt{list\_categories.php})}
\label{sec:orgd4d6507}
\begin{minted}[]{php}
<?php
require_once 'Categories.php';
require_once 'db_connect.php';

$categories = new Categories($pdo);
$allCategories = $categories->getAll();
?>

<!DOCTYPE html>
<html>
<head>
    <title>Category List</title>
</head>
<body>
    <h1>Category List</h1>
    <ul>
    <?php foreach ($allCategories as $category): ?>
        <li><?php echo htmlspecialchars($category['name']); ?></li>
    <?php endforeach; ?>
    </ul>
</body>
</html>
\end{minted}

\subsection{Security Considerations}
\label{sec:org4cc838c}
\begin{itemize}
\item \textbf{Input Validation}: Always validate and sanitize inputs to prevent SQL
injection and other forms of attacks.
\item \textbf{Error Handling}: Implement proper error handling, especially for
database operations.
\item \textbf{User Permissions}: Depending on your application, you might want to
restrict category management to certain user roles (like admin).
\end{itemize}

This class and the corresponding HTML forms provide basic functionality
for managing categories in our To-Do application. In a full
application, you'd also include methods for updating and deleting
categories, along with the necessary HTML forms and PHP scripts.
\end{document}
